% --- Archon physics formalization source begin ---
% archon:physics
% archon:covers PhyXMiniProblems/problem_phyx_mini_0959.lean
% archon:source-report reports/phyx_mini/problem_phyx_mini_0959.source.json

\chapter{Physics problem phyx\_mini\_0959}
\label{ch:PhyXMiniProblems_problem_phyx_mini_0959}

\paragraph{Problem source.}
\#\# Physical scenario

An electron and a proton are each moving at \textbackslash{}( 735\textbackslash{},\textbackslash{}text\{km/s\} \textbackslash{}) in perpendicular paths.

\#\# Question

At the instant when they are at the positions shown, the magnetic field the electron produces at the location of the proton.

\#\# Answer choices

- A: A: 1.12 \textbackslash{}times 10\textasciicircum{}\{-4\}\textbackslash{}T

- B: B: 2.24 \textbackslash{}times 10\textasciicircum{}\{-4\}\textbackslash{}T

- C: C: 2.24 \textbackslash{}times 10\textasciicircum{}\{-3\}\textbackslash{}T

- D: D: 3.36 \textbackslash{}times 10\textasciicircum{}\{-4\}\textbackslash{}T

\#\# Figure caption (auxiliary; use the image as primary evidence)

The image is a diagram with a coordinate system centered at point O. It shows two particles: an electron and a proton. 

- The electron is represented by a blue circle with a negative sign (-) inside it, located on the y-axis at the point (0, 5.00 nm). It has a green arrow pointing leftward, indicating direction or force.

- The proton is depicted as a red circle with a positive sign (+) inside it, located on the x-axis at the point (4.00 nm, 0). It has a green arrow pointing downward.

- The x and y axes are labeled, and the distances from the origin to each particle are marked as 4.00 nm for the proton and 5.00 nm for the electron. Additionally, the electron and proton are labeled on the diagram.

\#\# Dataset metadata

Magnetism; Physical Model Grounding Reasoning, Spatial Relation Reasoning

\paragraph{Recorded answer/context.}
B: B: 2.24 \textbackslash{}times 10\textasciicircum{}\{-4\}\textbackslash{}T

\paragraph{Figure/image path.}
/root/proposal\_for\_physic/science-mango/phyx\_mini\_run/phyx\_data/test\_image/959.png

\paragraph{Formalization target.}
create a compiling Lean file with sorry bodies at `PhyXMiniProblems/problem\_phyx\_mini\_0959.lean`. The Lean declarations must preserve the physical quantities, dimensions or dimensional roles, figure labels, governing-law hypotheses, and final relation expressed by this problem.
Use Mathlib/Physlib names found through LeanExplore where available. If a physics API is missing, introduce faithful local abstractions rather than scalar placeholder aliases.

\begin{theorem}[Physics formalization target]
\label{thm:physics:phyx_mini_0959:target}
\lean{PhyXMiniProblems.ProblemPhyXMini0959.problem_phyx_mini_0959}
\uses{def:physics:phyx-mini-0959:phyxminiproblems-problemphyxmini0959-movingelectronprotonsetup, def:physics:phyx-mini-0959:phyxminiproblems-problemphyxmini0959-electronfieldatprotoninteslas, def:physics:phyx-mini-0959:phyxminiproblems-problemphyxmini0959-electronfieldmagnitudeatprotoninteslas, def:physics:phyx-mini-0959:phyxminiproblems-problemphyxmini0959-matchesproblemstatementandprimaryfigure, def:physics:phyx-mini-0959:phyxminiproblems-problemphyxmini0959-hasphysicalmovingparticleparameters, def:physics:phyx-mini-0959:phyxminiproblems-problemphyxmini0959-usesstandardvacuumpermeability, def:physics:phyx-mini-0959:phyxminiproblems-problemphyxmini0959-satisfiesmovingpointchargemagneticfieldlaw, def:physics:phyx-mini-0959:phyxminiproblems-problemphyxmini0959-exactmodelfieldmagnitudeinteslas, def:physics:phyx-mini-0959:phyxminiproblems-problemphyxmini0959-pointsinaxisdirection, def:physics:phyx-mini-0959:phyxminiproblems-problemphyxmini0959-displayedmagneticfieldmagnitudeinteslas, def:physics:phyx-mini-0959:phyxminiproblems-problemphyxmini0959-recordedanswerchoice, def:physics:phyx-mini-0959:phyxminiproblems-problemphyxmini0959-roundstodisplayedprecision, def:physics:phyx-mini-0959:phyxminiproblems-problemphyxmini0959-isuniqueclosestdisplayedanswer}
The assigned autoformalize agent should translate this physics problem into Lean declarations in the covered file, with theorem and lemma proof bodies written as `by sorry`.
\end{theorem}
\begin{proof}
This is an autoformalization task, not a proof task. Produce statements that can later be proved by the physics prover without weakening the physical model.
\end{proof}

% --- Archon declaration topology begin ---
\section{Declaration topology}
\begin{definition}[Spatial Vector]
  \label{def:physics:phyx-mini-0959:phyxminiproblems-problemphyxmini0959-spatialvector}
  \lean{PhyXMiniProblems.ProblemPhyXMini0959.SpatialVector}
  Three-dimensional real vectors used only for coherent-unit readouts
\end{definition}
\begin{proof}
  This is the declared model definition of Spatial Vector.
\end{proof}
\begin{definition}[velocity Dimension]
  \label{def:physics:phyx-mini-0959:phyxminiproblems-problemphyxmini0959-velocitydimension}
  \lean{PhyXMiniProblems.ProblemPhyXMini0959.velocityDimension}
  The physical dimension L T⁻¹ of velocity
\end{definition}
\begin{proof}
  This is the declared model definition of velocity Dimension.
\end{proof}
\begin{definition}[magnetic Permeability Dimension]
  \label{def:physics:phyx-mini-0959:phyxminiproblems-problemphyxmini0959-magneticpermeabilitydimension}
  \lean{PhyXMiniProblems.ProblemPhyXMini0959.magneticPermeabilityDimension}
  Magnetic permeability has dimension M L C⁻²
\end{definition}
\begin{proof}
  This is the declared model definition of magnetic Permeability Dimension.
\end{proof}
\begin{definition}[Position Vector Quantity]
  \label{def:physics:phyx-mini-0959:phyxminiproblems-problemphyxmini0959-positionvectorquantity}
  \lean{PhyXMiniProblems.ProblemPhyXMini0959.PositionVectorQuantity}
  \uses{def:physics:phyx-mini-0959:phyxminiproblems-problemphyxmini0959-spatialvector}
  A unit-independent signed three-dimensional position vector
\end{definition}
\begin{proof}
  This is the declared model definition of Position Vector Quantity.
\end{proof}
\begin{definition}[Velocity Vector Quantity]
  \label{def:physics:phyx-mini-0959:phyxminiproblems-problemphyxmini0959-velocityvectorquantity}
  \lean{PhyXMiniProblems.ProblemPhyXMini0959.VelocityVectorQuantity}
  \uses{def:physics:phyx-mini-0959:phyxminiproblems-problemphyxmini0959-spatialvector, def:physics:phyx-mini-0959:phyxminiproblems-problemphyxmini0959-velocitydimension}
  A unit-independent signed three-dimensional velocity vector
\end{definition}
\begin{proof}
  This is the declared model definition of Velocity Vector Quantity.
\end{proof}
\begin{definition}[Signed Charge Quantity]
  \label{def:physics:phyx-mini-0959:phyxminiproblems-problemphyxmini0959-signedchargequantity}
  \lean{PhyXMiniProblems.ProblemPhyXMini0959.SignedChargeQuantity}
  A signed, unit-independent electric charge
\end{definition}
\begin{proof}
  This is the declared model definition of Signed Charge Quantity.
\end{proof}
\begin{definition}[Magnetic Permeability Quantity]
  \label{def:physics:phyx-mini-0959:phyxminiproblems-problemphyxmini0959-magneticpermeabilityquantity}
  \lean{PhyXMiniProblems.ProblemPhyXMini0959.MagneticPermeabilityQuantity}
  \uses{def:physics:phyx-mini-0959:phyxminiproblems-problemphyxmini0959-magneticpermeabilitydimension}
  A nonnegative, unit-independent magnetic permeability
\end{definition}
\begin{proof}
  This is the declared model definition of Magnetic Permeability Quantity.
\end{proof}
\begin{definition}[position Vector Readout]
  \label{def:physics:phyx-mini-0959:phyxminiproblems-problemphyxmini0959-positionvectorreadout}
  \lean{PhyXMiniProblems.ProblemPhyXMini0959.positionVectorReadout}
  \uses{def:physics:phyx-mini-0959:phyxminiproblems-problemphyxmini0959-spatialvector, def:physics:phyx-mini-0959:phyxminiproblems-problemphyxmini0959-positionvectorquantity}
  Read a physical position vector in a selected coherent length unit
\end{definition}
\begin{proof}
  This is the declared model definition of position Vector Readout.
\end{proof}
\begin{definition}[position Vector In Meters]
  \label{def:physics:phyx-mini-0959:phyxminiproblems-problemphyxmini0959-positionvectorinmeters}
  \lean{PhyXMiniProblems.ProblemPhyXMini0959.positionVectorInMeters}
  \uses{def:physics:phyx-mini-0959:phyxminiproblems-problemphyxmini0959-spatialvector, def:physics:phyx-mini-0959:phyxminiproblems-problemphyxmini0959-positionvectorquantity, def:physics:phyx-mini-0959:phyxminiproblems-problemphyxmini0959-positionvectorreadout}
  Read a physical position vector in coherent-SI metres
\end{definition}
\begin{proof}
  This is the declared model definition of position Vector In Meters.
\end{proof}
\begin{definition}[position Vector In Nanometers]
  \label{def:physics:phyx-mini-0959:phyxminiproblems-problemphyxmini0959-positionvectorinnanometers}
  \lean{PhyXMiniProblems.ProblemPhyXMini0959.positionVectorInNanometers}
  \uses{def:physics:phyx-mini-0959:phyxminiproblems-problemphyxmini0959-spatialvector, def:physics:phyx-mini-0959:phyxminiproblems-problemphyxmini0959-positionvectorquantity, def:physics:phyx-mini-0959:phyxminiproblems-problemphyxmini0959-positionvectorreadout}
  Read a physical position vector in nanometres, as labelled in the raster
\end{definition}
\begin{proof}
  This is the declared model definition of position Vector In Nanometers.
\end{proof}
\begin{definition}[velocity Vector In Meters Per Second]
  \label{def:physics:phyx-mini-0959:phyxminiproblems-problemphyxmini0959-velocityvectorinmeterspersecond}
  \lean{PhyXMiniProblems.ProblemPhyXMini0959.velocityVectorInMetersPerSecond}
  \uses{def:physics:phyx-mini-0959:phyxminiproblems-problemphyxmini0959-spatialvector, def:physics:phyx-mini-0959:phyxminiproblems-problemphyxmini0959-velocityvectorquantity}
  Read a physical velocity vector in coherent-SI metres per second
\end{definition}
\begin{proof}
  This is the declared model definition of velocity Vector In Meters Per Second.
\end{proof}
\begin{definition}[charge In Coulombs]
  \label{def:physics:phyx-mini-0959:phyxminiproblems-problemphyxmini0959-chargeincoulombs}
  \lean{PhyXMiniProblems.ProblemPhyXMini0959.chargeInCoulombs}
  \uses{def:physics:phyx-mini-0959:phyxminiproblems-problemphyxmini0959-signedchargequantity}
  Read a signed charge in coherent-SI coulombs
\end{definition}
\begin{proof}
  This is the declared model definition of charge In Coulombs.
\end{proof}
\begin{definition}[charge In Elementary Charges]
  \label{def:physics:phyx-mini-0959:phyxminiproblems-problemphyxmini0959-chargeinelementarycharges}
  \lean{PhyXMiniProblems.ProblemPhyXMini0959.chargeInElementaryCharges}
  \uses{def:physics:phyx-mini-0959:phyxminiproblems-problemphyxmini0959-signedchargequantity}
  Read a signed charge in Physlib elementary-charge units
\end{definition}
\begin{proof}
  This is the declared model definition of charge In Elementary Charges.
\end{proof}
\begin{definition}[speed Readout]
  \label{def:physics:phyx-mini-0959:phyxminiproblems-problemphyxmini0959-speedreadout}
  \lean{PhyXMiniProblems.ProblemPhyXMini0959.speedReadout}
  Read a speed in a coherent choice of length and time units
\end{definition}
\begin{proof}
  This is the declared model definition of speed Readout.
\end{proof}
\begin{definition}[permeability In Tesla Meters Per Ampere]
  \label{def:physics:phyx-mini-0959:phyxminiproblems-problemphyxmini0959-permeabilityinteslametersperampere}
  \lean{PhyXMiniProblems.ProblemPhyXMini0959.permeabilityInTeslaMetersPerAmpere}
  \uses{def:physics:phyx-mini-0959:phyxminiproblems-problemphyxmini0959-magneticpermeabilityquantity}
  Read magnetic permeability in coherent-SI tesla-metres per ampere
\end{definition}
\begin{proof}
  This is the declared model definition of permeability In Tesla Meters Per Ampere.
\end{proof}
\begin{definition}[spatial Cross Product]
  \label{def:physics:phyx-mini-0959:phyxminiproblems-problemphyxmini0959-spatialcrossproduct}
  \lean{PhyXMiniProblems.ProblemPhyXMini0959.spatialCrossProduct}
  \uses{def:physics:phyx-mini-0959:phyxminiproblems-problemphyxmini0959-spatialvector}
  Mathlib's coordinate cross product transported to Euclidean space
\end{definition}
\begin{proof}
  This is the declared model definition of spatial Cross Product.
\end{proof}
\begin{definition}[Coordinate Axis]
  \label{def:physics:phyx-mini-0959:phyxminiproblems-problemphyxmini0959-coordinateaxis}
  \lean{PhyXMiniProblems.ProblemPhyXMini0959.CoordinateAxis}
  The three coordinate axes, with z normal to the displayed page
\end{definition}
\begin{proof}
  This is the declared model definition of Coordinate Axis.
\end{proof}
\begin{definition}[axis Vector]
  \label{def:physics:phyx-mini-0959:phyxminiproblems-problemphyxmini0959-axisvector}
  \lean{PhyXMiniProblems.ProblemPhyXMini0959.axisVector}
  \uses{def:physics:phyx-mini-0959:phyxminiproblems-problemphyxmini0959-spatialvector, def:physics:phyx-mini-0959:phyxminiproblems-problemphyxmini0959-coordinateaxis}
  Standard positive unit vector along a coordinate axis
\end{definition}
\begin{proof}
  This is the declared model definition of axis Vector.
\end{proof}
\begin{definition}[Particle]
  \label{def:physics:phyx-mini-0959:phyxminiproblems-problemphyxmini0959-particle}
  \lean{PhyXMiniProblems.ProblemPhyXMini0959.Particle}
  The two particles explicitly named in the problem and raster
\end{definition}
\begin{proof}
  This is the declared model definition of Particle.
\end{proof}
\begin{definition}[Charge Sign Glyph]
  \label{def:physics:phyx-mini-0959:phyxminiproblems-problemphyxmini0959-chargesignglyph}
  \lean{PhyXMiniProblems.ProblemPhyXMini0959.ChargeSignGlyph}
  Charge-sign glyph drawn inside a particle marker
\end{definition}
\begin{proof}
  This is the declared model definition of Charge Sign Glyph.
\end{proof}
\begin{definition}[Particle Fill Color]
  \label{def:physics:phyx-mini-0959:phyxminiproblems-problemphyxmini0959-particlefillcolor}
  \lean{PhyXMiniProblems.ProblemPhyXMini0959.ParticleFillColor}
  Fill colours visibly distinguishing the two particle markers
\end{definition}
\begin{proof}
  This is the declared model definition of Particle Fill Color.
\end{proof}
\begin{definition}[Axis Direction]
  \label{def:physics:phyx-mini-0959:phyxminiproblems-problemphyxmini0959-axisdirection}
  \lean{PhyXMiniProblems.ProblemPhyXMini0959.AxisDirection}
  Oriented Cartesian directions used by the two velocity arrows
\end{definition}
\begin{proof}
  This is the declared model definition of Axis Direction.
\end{proof}
\begin{definition}[direction Vector]
  \label{def:physics:phyx-mini-0959:phyxminiproblems-problemphyxmini0959-directionvector}
  \lean{PhyXMiniProblems.ProblemPhyXMini0959.directionVector}
  \uses{def:physics:phyx-mini-0959:phyxminiproblems-problemphyxmini0959-spatialvector, def:physics:phyx-mini-0959:phyxminiproblems-problemphyxmini0959-axisvector, def:physics:phyx-mini-0959:phyxminiproblems-problemphyxmini0959-axisdirection}
  Coherent-SI unit vector associated with an oriented direction
\end{definition}
\begin{proof}
  This is the declared model definition of direction Vector.
\end{proof}
\begin{definition}[Electron Proton Figure]
  \label{def:physics:phyx-mini-0959:phyxminiproblems-problemphyxmini0959-electronprotonfigure}
  \lean{PhyXMiniProblems.ProblemPhyXMini0959.ElectronProtonFigure}
  \uses{def:physics:phyx-mini-0959:phyxminiproblems-problemphyxmini0959-coordinateaxis, def:physics:phyx-mini-0959:phyxminiproblems-problemphyxmini0959-particle, def:physics:phyx-mini-0959:phyxminiproblems-problemphyxmini0959-chargesignglyph, def:physics:phyx-mini-0959:phyxminiproblems-problemphyxmini0959-particlefillcolor, def:physics:phyx-mini-0959:phyxminiproblems-problemphyxmini0959-axisdirection}
  Literal content transcribed from the supplied raster 959.png
\end{definition}
\begin{proof}
  This is the declared model definition of Electron Proton Figure.
\end{proof}
\begin{definition}[Moving Electron Proton Setup]
  \label{def:physics:phyx-mini-0959:phyxminiproblems-problemphyxmini0959-movingelectronprotonsetup}
  \lean{PhyXMiniProblems.ProblemPhyXMini0959.MovingElectronProtonSetup}
  \uses{def:physics:phyx-mini-0959:phyxminiproblems-problemphyxmini0959-positionvectorquantity, def:physics:phyx-mini-0959:phyxminiproblems-problemphyxmini0959-velocityvectorquantity, def:physics:phyx-mini-0959:phyxminiproblems-problemphyxmini0959-signedchargequantity, def:physics:phyx-mini-0959:phyxminiproblems-problemphyxmini0959-magneticpermeabilityquantity, def:physics:phyx-mini-0959:phyxminiproblems-problemphyxmini0959-particle, def:physics:phyx-mini-0959:phyxminiproblems-problemphyxmini0959-electronprotonfigure}
  Independent physical objects at the pictured instant. In particular, the electron's magnetic field is not defined from the displayed answer
\end{definition}
\begin{proof}
  This is the declared model definition of Moving Electron Proton Setup.
\end{proof}
\begin{definition}[electron To Proton Displacement In Meters]
  \label{def:physics:phyx-mini-0959:phyxminiproblems-problemphyxmini0959-electrontoprotondisplacementinmeters}
  \lean{PhyXMiniProblems.ProblemPhyXMini0959.electronToProtonDisplacementInMeters}
  \uses{def:physics:phyx-mini-0959:phyxminiproblems-problemphyxmini0959-spatialvector, def:physics:phyx-mini-0959:phyxminiproblems-problemphyxmini0959-positionvectorinmeters, def:physics:phyx-mini-0959:phyxminiproblems-problemphyxmini0959-movingelectronprotonsetup}
  Displacement from the electron to the proton, read in metres
\end{definition}
\begin{proof}
  This is the declared model definition of electron To Proton Displacement In Meters.
\end{proof}
\begin{definition}[space Point Of Meters]
  \label{def:physics:phyx-mini-0959:phyxminiproblems-problemphyxmini0959-spacepointofmeters}
  \lean{PhyXMiniProblems.ProblemPhyXMini0959.spacePointOfMeters}
  \uses{def:physics:phyx-mini-0959:phyxminiproblems-problemphyxmini0959-spatialvector}
  Interpret a coherent-SI metre vector as a Physlib spatial point
\end{definition}
\begin{proof}
  This is the declared model definition of space Point Of Meters.
\end{proof}
\begin{definition}[electron Field At Proton In Teslas]
  \label{def:physics:phyx-mini-0959:phyxminiproblems-problemphyxmini0959-electronfieldatprotoninteslas}
  \lean{PhyXMiniProblems.ProblemPhyXMini0959.electronFieldAtProtonInTeslas}
  \uses{def:physics:phyx-mini-0959:phyxminiproblems-problemphyxmini0959-spatialvector, def:physics:phyx-mini-0959:phyxminiproblems-problemphyxmini0959-positionvectorinmeters, def:physics:phyx-mini-0959:phyxminiproblems-problemphyxmini0959-movingelectronprotonsetup, def:physics:phyx-mini-0959:phyxminiproblems-problemphyxmini0959-spacepointofmeters}
  Electron-produced magnetic-field vector at the proton, in teslas
\end{definition}
\begin{proof}
  This is the declared model definition of electron Field At Proton In Teslas.
\end{proof}
\begin{definition}[electron Field Magnitude At Proton In Teslas]
  \label{def:physics:phyx-mini-0959:phyxminiproblems-problemphyxmini0959-electronfieldmagnitudeatprotoninteslas}
  \lean{PhyXMiniProblems.ProblemPhyXMini0959.electronFieldMagnitudeAtProtonInTeslas}
  \uses{def:physics:phyx-mini-0959:phyxminiproblems-problemphyxmini0959-movingelectronprotonsetup, def:physics:phyx-mini-0959:phyxminiproblems-problemphyxmini0959-electronfieldatprotoninteslas}
  Magnitude of the electron-produced field at the proton, in teslas
\end{definition}
\begin{proof}
  This is the declared model definition of electron Field Magnitude At Proton In Teslas.
\end{proof}
\begin{definition}[Matches Problem Statement And Primary Figure]
  \label{def:physics:phyx-mini-0959:phyxminiproblems-problemphyxmini0959-matchesproblemstatementandprimaryfigure}
  \lean{PhyXMiniProblems.ProblemPhyXMini0959.MatchesProblemStatementAndPrimaryFigure}
  \uses{def:physics:phyx-mini-0959:phyxminiproblems-problemphyxmini0959-positionvectorinmeters, def:physics:phyx-mini-0959:phyxminiproblems-problemphyxmini0959-positionvectorinnanometers, def:physics:phyx-mini-0959:phyxminiproblems-problemphyxmini0959-velocityvectorinmeterspersecond, def:physics:phyx-mini-0959:phyxminiproblems-problemphyxmini0959-chargeincoulombs, def:physics:phyx-mini-0959:phyxminiproblems-problemphyxmini0959-chargeinelementarycharges, def:physics:phyx-mini-0959:phyxminiproblems-problemphyxmini0959-speedreadout, def:physics:phyx-mini-0959:phyxminiproblems-problemphyxmini0959-axisvector, def:physics:phyx-mini-0959:phyxminiproblems-problemphyxmini0959-directionvector, def:physics:phyx-mini-0959:phyxminiproblems-problemphyxmini0959-movingelectronprotonsetup}
  Written data and all literal/calibrated evidence from the primary image
\end{definition}
\begin{proof}
  This is the declared model definition of Matches Problem Statement And Primary Figure.
\end{proof}
\begin{definition}[Has Physical Moving Particle Parameters]
  \label{def:physics:phyx-mini-0959:phyxminiproblems-problemphyxmini0959-hasphysicalmovingparticleparameters}
  \lean{PhyXMiniProblems.ProblemPhyXMini0959.HasPhysicalMovingParticleParameters}
  \uses{def:physics:phyx-mini-0959:phyxminiproblems-problemphyxmini0959-velocityvectorinmeterspersecond, def:physics:phyx-mini-0959:phyxminiproblems-problemphyxmini0959-chargeincoulombs, def:physics:phyx-mini-0959:phyxminiproblems-problemphyxmini0959-speedreadout, def:physics:phyx-mini-0959:phyxminiproblems-problemphyxmini0959-permeabilityinteslametersperampere, def:physics:phyx-mini-0959:phyxminiproblems-problemphyxmini0959-movingelectronprotonsetup, def:physics:phyx-mini-0959:phyxminiproblems-problemphyxmini0959-electrontoprotondisplacementinmeters}
  Positivity and nondegeneracy conditions selecting the physical branch
\end{definition}
\begin{proof}
  This is the declared model definition of Has Physical Moving Particle Parameters.
\end{proof}
\begin{definition}[Uses Standard Vacuum Permeability]
  \label{def:physics:phyx-mini-0959:phyxminiproblems-problemphyxmini0959-usesstandardvacuumpermeability}
  \lean{PhyXMiniProblems.ProblemPhyXMini0959.UsesStandardVacuumPermeability}
  \uses{def:physics:phyx-mini-0959:phyxminiproblems-problemphyxmini0959-permeabilityinteslametersperampere, def:physics:phyx-mini-0959:phyxminiproblems-problemphyxmini0959-movingelectronprotonsetup}
  The dimensionful vacuum permeability agrees with Physlib's electromagnetic system and has the standard SI value 4π * 10⁻⁷ T m/A
\end{definition}
\begin{proof}
  This is the declared model definition of Uses Standard Vacuum Permeability.
\end{proof}
\begin{definition}[Satisfies Moving Point Charge Magnetic Field Law]
  \label{def:physics:phyx-mini-0959:phyxminiproblems-problemphyxmini0959-satisfiesmovingpointchargemagneticfieldlaw}
  \lean{PhyXMiniProblems.ProblemPhyXMini0959.SatisfiesMovingPointChargeMagneticFieldLaw}
  \uses{def:physics:phyx-mini-0959:phyxminiproblems-problemphyxmini0959-velocityvectorinmeterspersecond, def:physics:phyx-mini-0959:phyxminiproblems-problemphyxmini0959-chargeincoulombs, def:physics:phyx-mini-0959:phyxminiproblems-problemphyxmini0959-permeabilityinteslametersperampere, def:physics:phyx-mini-0959:phyxminiproblems-problemphyxmini0959-spatialcrossproduct, def:physics:phyx-mini-0959:phyxminiproblems-problemphyxmini0959-movingelectronprotonsetup, def:physics:phyx-mini-0959:phyxminiproblems-problemphyxmini0959-electrontoprotondisplacementinmeters, def:physics:phyx-mini-0959:phyxminiproblems-problemphyxmini0959-electronfieldatprotoninteslas}
  Low-speed field of a moving point charge at the proton: B = μ₀ q /(4π |r|³) (v × r). Here r points from the electron to the proton. This is a governing law in the independent physical quantities and contains no evaluated field or answer choice
\end{definition}
\begin{proof}
  This is the declared model definition of Satisfies Moving Point Charge Magnetic Field Law.
\end{proof}
\begin{definition}[exact Model Field Magnitude In Teslas]
  \label{def:physics:phyx-mini-0959:phyxminiproblems-problemphyxmini0959-exactmodelfieldmagnitudeinteslas}
  \lean{PhyXMiniProblems.ProblemPhyXMini0959.exactModelFieldMagnitudeInTeslas}
  Exact positive field magnitude produced by the model, in teslas
\end{definition}
\begin{proof}
  This is the declared model definition of exact Model Field Magnitude In Teslas.
\end{proof}
\begin{definition}[Points In Axis Direction]
  \label{def:physics:phyx-mini-0959:phyxminiproblems-problemphyxmini0959-pointsinaxisdirection}
  \lean{PhyXMiniProblems.ProblemPhyXMini0959.PointsInAxisDirection}
  \uses{def:physics:phyx-mini-0959:phyxminiproblems-problemphyxmini0959-spatialvector, def:physics:phyx-mini-0959:phyxminiproblems-problemphyxmini0959-axisdirection, def:physics:phyx-mini-0959:phyxminiproblems-problemphyxmini0959-directionvector}
  A nonzero field vector points in a specified oriented axis direction
\end{definition}
\begin{proof}
  This is the declared model definition of Points In Axis Direction.
\end{proof}
\begin{lemma}[electron Field At Proton exact]
  \label{lem:physics:phyx-mini-0959:phyxminiproblems-problemphyxmini0959-electronfieldatproton-exact}
  \lean{PhyXMiniProblems.ProblemPhyXMini0959.electronFieldAtProton_exact}
  \uses{def:physics:phyx-mini-0959:phyxminiproblems-problemphyxmini0959-directionvector, def:physics:phyx-mini-0959:phyxminiproblems-problemphyxmini0959-movingelectronprotonsetup, def:physics:phyx-mini-0959:phyxminiproblems-problemphyxmini0959-electronfieldatprotoninteslas, def:physics:phyx-mini-0959:phyxminiproblems-problemphyxmini0959-matchesproblemstatementandprimaryfigure, def:physics:phyx-mini-0959:phyxminiproblems-problemphyxmini0959-hasphysicalmovingparticleparameters, def:physics:phyx-mini-0959:phyxminiproblems-problemphyxmini0959-usesstandardvacuumpermeability, def:physics:phyx-mini-0959:phyxminiproblems-problemphyxmini0959-satisfiesmovingpointchargemagneticfieldlaw, def:physics:phyx-mini-0959:phyxminiproblems-problemphyxmini0959-exactmodelfieldmagnitudeinteslas}
  The moving-charge law and pictured geometry give an into-page (-z) field with the exact model magnitude above
\end{lemma}
\begin{proof}
  The conclusion follows from the governing relations and figure readouts named in the statement of electron Field At Proton exact.
\end{proof}
\begin{definition}[Answer Choice]
  \label{def:physics:phyx-mini-0959:phyxminiproblems-problemphyxmini0959-answerchoice}
  \lean{PhyXMiniProblems.ProblemPhyXMini0959.AnswerChoice}
  Labels of the four magnetic-field magnitudes printed in the question
\end{definition}
\begin{proof}
  This is the declared model definition of Answer Choice.
\end{proof}
\begin{definition}[displayed Magnetic Field Magnitude In Teslas]
  \label{def:physics:phyx-mini-0959:phyxminiproblems-problemphyxmini0959-displayedmagneticfieldmagnitudeinteslas}
  \lean{PhyXMiniProblems.ProblemPhyXMini0959.displayedMagneticFieldMagnitudeInTeslas}
  \uses{def:physics:phyx-mini-0959:phyxminiproblems-problemphyxmini0959-answerchoice}
  Magnetic-field magnitude in teslas displayed beside each choice
\end{definition}
\begin{proof}
  This is the declared model definition of displayed Magnetic Field Magnitude In Teslas.
\end{proof}
\begin{definition}[recorded Answer Choice]
  \label{def:physics:phyx-mini-0959:phyxminiproblems-problemphyxmini0959-recordedanswerchoice}
  \lean{PhyXMiniProblems.ProblemPhyXMini0959.recordedAnswerChoice}
  \uses{def:physics:phyx-mini-0959:phyxminiproblems-problemphyxmini0959-answerchoice}
  The answer label recorded by the source dataset
\end{definition}
\begin{proof}
  This is the declared model definition of recorded Answer Choice.
\end{proof}
\begin{definition}[Rounds To Displayed Precision]
  \label{def:physics:phyx-mini-0959:phyxminiproblems-problemphyxmini0959-roundstodisplayedprecision}
  \lean{PhyXMiniProblems.ProblemPhyXMini0959.RoundsToDisplayedPrecision}
  Rounding to the displayed 0.01 * 10⁻⁴ T precision
\end{definition}
\begin{proof}
  This is the declared model definition of Rounds To Displayed Precision.
\end{proof}
\begin{definition}[Is Unique Closest Displayed Answer]
  \label{def:physics:phyx-mini-0959:phyxminiproblems-problemphyxmini0959-isuniqueclosestdisplayedanswer}
  \lean{PhyXMiniProblems.ProblemPhyXMini0959.IsUniqueClosestDisplayedAnswer}
  \uses{def:physics:phyx-mini-0959:phyxminiproblems-problemphyxmini0959-answerchoice, def:physics:phyx-mini-0959:phyxminiproblems-problemphyxmini0959-displayedmagneticfieldmagnitudeinteslas}
  A choice is strictly closer than every distinct displayed alternative
\end{definition}
\begin{proof}
  This is the declared model definition of Is Unique Closest Displayed Answer.
\end{proof}
% --- Archon declaration topology end ---
% --- Archon physics formalization source end ---
